% This file was converted to LaTeX by Writer2LaTeX ver. 1.9.9
% see http://writer2latex.sourceforge.net for more info
\documentclass{article}
\usepackage{calc,amsmath,amssymb,amsfonts}
\usepackage[LGR,T1]{fontenc}
\usepackage[greek,french]{babel}
\usepackage[style=numeric,backend=biber]{biblatex}
\title{Performance et confiance : un leadership qui commence par le respect}
\author{David Graz}
\date{2026-04-01}
\begin{document}
\clearpage\title{Performance et confiance : un leadership qui commence par le respect}
\maketitle

Graz Network


\bigskip


\bigskip


\bigskip


\bigskip

Résumé : Cet article soutient une affirmation claire\leavevmode\allowbreak\,: la performance durable ne se
$\text{\textgreek{«}}$\leavevmode\allowbreak\,décrète\leavevmode\allowbreak\,$\text{\textgreek{»}}$ pas, elle émerge
dans un contexte qui encourage l$\text{\textgreek{’}}$engagement, l$\text{\textgreek{’}}$initiative et
l$\text{\textgreek{’}}$apprentissage. Ce contexte s$\text{\textgreek{’}}$appelle la confiance, et la confiance ne se
construit pas sans une condition préalable\leavevmode\allowbreak\,: le respect. En s$\text{\textgreek{’}}$appuyant sur
une analyse $\text{\textgreek{«}}$\leavevmode\allowbreak\,intégrative\leavevmode\allowbreak\,$\text{\textgreek{»}}$ du
leadership, il est important de considérer les traits (dispositions individuelles) des leaders, mais ce sont surtout
les comportements observables qui structurent l$\text{\textgreek{’}}$environnement de travail et qui expliquent une
grande partie de l$\text{\textgreek{’}}$efficacité collective. En conséquence, il est possible
d$\text{\textgreek{’}}$entraîner et d$\text{\textgreek{’}}$installer des comportements qui rendent le respect concret
et la confiance plausible. Cette logique est cohérente avec les recherche montrant que la confiance médie une partie du
lien leadership–performance \label{ref:RNDxgeyzVUFHQ}(Dirks \& Ferrin, 2002; Legood et al., 2021) et que la sécurité
psychologique est un mécanisme central reliant comportements de leadership, apprentissage d$\text{\textgreek{’}}$équipe
et résultats \label{ref:RNDo3M3FyKsg2}(Frazier et al., 2017; Koeslag-Kreunen et al., 2018).

\clearpage
La performance fascine, mesure, compare et sélectionne. Pourtant, elle ne se décrète pas : elle se libère. Et dans une
équipe, un collectif ou une organisation, la performance durable apparaît rarement comme un simple produit
d$\text{\textgreek{’}}$exigences. Généralement, elle émerge d$\text{\textgreek{’}}$un climat qui autorise
l$\text{\textgreek{’}}$engagement, l$\text{\textgreek{’}}$initiative et l$\text{\textgreek{’}}$apprentissage. Ce climat
porte un nom : la confiance.

Or, la confiance ne naît pas d$\text{\textgreek{’}}$un discours. Elle se construit. Et les chemins qui y conduisent
passent presque toujours par une condition première : le respect. C$\text{\textgreek{’}}$est ainsi que la performance
peut s$\text{\textgreek{’}}$inscrire dans un enchaînement simple et exigeant : respect, confiance, performance. 

Une telle pensée s$\text{\textgreek{’}}$accorde alors avec une lecture intégrative du leadership. Les dispositions
individuelles (traits) influencent les comportements observables, qui constituent un mécanisme central expliquant
l$\text{\textgreek{’}}$efficacité et les résultats collectifs. \label{ref:RNDs9ofFbAlk6}(Derue et al., 2011). 

\section{Construire un environnement}
Selon Kant, la moralité se juge par le devoir et l$\text{\textgreek{’}}$intention. Et ensuite, par
l$\text{\textgreek{’}}$existence de principes suffisamment robustes pour être appliqués de manière cohérente. La raison
fonde alors des règles universelles \label{ref:RNDkGpEPQ9nT6}(Kant et al., 2008). Un leadership qui pose le cadre peut
donc être lu comme une traduction pratique qui stabilise des normes tenables par tous. 

Cela rejoint l$\text{\textgreek{’}}$impératif d$\text{\textgreek{’}}$agir selon une maxime fondée sur le respect. Car il
protège le cadre, rend la relation lisible, et responsabilise l$\text{\textgreek{’}}$autorité. Et, dans une lecture
$\text{\textgreek{«}}$ développementale $\text{\textgreek{»}}$ du leadership, un environnement favorable
n$\text{\textgreek{’}}$est pas un bonus : c$\text{\textgreek{’}}$est une condition pour que les apprentissages et les
comportements souhaités puissent réellement s$\text{\textgreek{’}}$ancrer \label{ref:RNDdzDKSYVG9u}(Wallace et al.,
2021). 

Stabiliser le respect revient d$\text{\textgreek{’}}$abord à stabiliser un cadre. Sur un plan déontologique,
$\text{\textgreek{«}}$ poser un cadre $\text{\textgreek{»}}$ (règles, rituels, recadrages) n$\text{\textgreek{’}}$est
donc pas une simple préférence managériale. C$\text{\textgreek{’}}$est une traduction opérationnelle du principe selon
lequel un collectif fonctionne mieux lorsque les normes sont tenables par tous, car elles réduisent
l$\text{\textgreek{’}}$arbitraire et rendent les attentes accessibles.

\subsubsection{Principes de leadership associés}
\begin{itemize}
\item Clarté du cadre : la clarté des règles réduit l$\text{\textgreek{’}}$arbitraire.
\item Cohérence : le cadre n$\text{\textgreek{’}}$a de valeur que s$\text{\textgreek{’}}$il est tenu.
\item Justice relationnelle : condamner par le dialogue.
\end{itemize}
Dans cette perspective $\text{\textgreek{«}}$le respect est une base$\text{\textgreek{»}}$, parce
qu$\text{\textgreek{’}}$il protège la relation contre l$\text{\textgreek{’}}$approximation et
l$\text{\textgreek{’}}$arbitraire. Un cadre respectueux n$\text{\textgreek{’}}$a toutefois de valeur que
s$\text{\textgreek{’}}$il est cohérent. Il doit être tenu (surtout lorsque cela coûte) et entretenu. 

Concrètement, ce cadre peut se matérialiser par des rituels où chacun explicite son objectif et son intention, où
l$\text{\textgreek{’}}$effort et la progression sont reconnus. Mais également des codes de conduite simples et
accessibles au plus grand nombre. Un rappel visible et transparent des standards est essentiel. 

Par ailleurs, lorsqu$\text{\textgreek{’}}$un écart survient, l$\text{\textgreek{’}}$exigence de respect impose un
recadrage par le dialogue : comprendre, expliciter l$\text{\textgreek{’}}$impact, rappeler la règle, exiger une
réparation proportionnée et éducative. Cette logique rend respectable l$\text{\textgreek{’}}$autorité et transforme la
sanction en acte de responsabilité plutôt qu$\text{\textgreek{’}}$en geste de domination.

\subsubsection{Actions concrètes}
\begin{itemize}
\item Rituels quotidiens : ouverture par un partage (objectif individuel + intention collective).
\item Rituels de clôture : clôture orientée reconnaissance (effort, progrès, entraide).
\item Code de conduite : simple, public, applicable, avec conséquences explicites.
\item Affichage des règles : rappel permanent des standards.
\item Sanctionner les écarts par le dialogue : préserver la dignité tout en protégeant le cadre.
\item Conséquences éducatives : sanctions utiles (réparation, contribution, apprentissage). 
\end{itemize}
Un tel socle est une protection contre certains excès personnels. Les travaux sur les $\text{\textgreek{«}}$ traits
$\text{\textgreek{»}}$ rappellent que les attributs de personnalité ne sont pas $\text{\textgreek{«}}$ bons
$\text{\textgreek{»}}$ ou $\text{\textgreek{«}}$ mauvais $\text{\textgreek{»}}$ en soi. Ils peuvent présenter des faces
lumineuses et sombres, selon la manière dont ils s$\text{\textgreek{’}}$expriment et selon les contextes. Un leadership
très consciencieux peut basculer en perfectionnisme coûteux si le cadre n$\text{\textgreek{’}}$impose pas des limites
de temps, de priorité et de proportionnalité. Comprendre ses inclinations (traits) permet d$\text{\textgreek{’}}$éviter
que ses forces deviennent des vulnérabilités \label{ref:RND3eZlw8jUor}(Judge et al., 2009). 

Ainsi, le respect n$\text{\textgreek{’}}$est pas seulement relationnel, il sert aussi de garde-fou contre des dérives.
Il \ augmente la crédibilité du leadership instauré. Le respect du cadre devient alors une hygiène de relationnelle. Il
empêche qu$\text{\textgreek{’}}$une force se transforme en faiblesse et renforce la crédibilité morale du leadership
exercé.

\section{Développer la confiance}
La confiance constitue à la fois un objectif relationnel et un levier d$\text{\textgreek{’}}$apprentissage. Elle ne
dépend pas uniquement du résultat, mais elle dépend de la qualité du lien et de la perception
d$\text{\textgreek{’}}$une intention justifiable. Se nourrissant de constance, de transparence, et
d$\text{\textgreek{’}}$un soutien crédible. Les synthèses de la recherche soulignent justement que les théories
contemporaines accordent une place majeure au contexte social et relationnel. Elles décrivent le leadership comme un
phénomène multiforme, sensible aux dynamiques de relations et aux environnements d$\text{\textgreek{’}}$action
\label{ref:RNDrhGpnvSF8o}(Dinh et al., 2014). Dans cette logique, la confiance n$\text{\textgreek{’}}$est pas un plus :
elle est un état collectif qui conditionne l$\text{\textgreek{’}}$initiative, la coopération et
l$\text{\textgreek{’}}$apprentissage. Ainsi, dans une approche moderne de développement, cela relève autant du leader
development (capacités individuelles) que du leadership development (états collectifs). Cette évolution se joue dans
l$\text{\textgreek{’}}$expérience : sortir de la zone de confort via des missions challengeantes (complexité, enjeux,
pression, nouveauté) fait grandir, à condition d$\text{\textgreek{’}}$être soutenu et de donner du sens
\label{ref:RNDxkLVNbumpB}(Day et al., 2021).

\subsubsection{Principes de leadership associés}
\begin{itemize}
\item Soutien sans complaisance : croire en l$\text{\textgreek{’}}$autre en le confrontant au réel.
\item Feedback orienté croissance : apprendre sans abîmer l$\text{\textgreek{’}}$estime.
\item Disponibilité et fiabilité : tenue des engagements.
\end{itemize}
La confiance se construit par une série d$\text{\textgreek{’}}$actes répétés. Ainsi, le feedback positif régulier
renforce l$\text{\textgreek{’}}$estime et la sécurité d$\text{\textgreek{’}}$oser. La critique constructive soutient
l$\text{\textgreek{’}}$amélioration sans abîmer l$\text{\textgreek{’}}$identité. Le suivi individualisé montre que le
potentiel de chacun est pris au sérieux. La coconstruction d$\text{\textgreek{’}}$objectifs mesurables et
d$\text{\textgreek{’}}$étapes concrètes transforme la confiance en contrat de progression. Ici, la transparence joue un
rôle décisif : expliquer les décisions, assumer les arbitrages, partager les erreurs et les leçons apprises. Cela rend
l$\text{\textgreek{’}}$intention lisible et l$\text{\textgreek{’}}$action crédible. Enfin, la disponibilité dans les
moments difficiles est la forme la plus coûteuse et donc la plus probante de l$\text{\textgreek{’}}$engagement. La
confiance se mesure moins à ce qui est déclaré qu$\text{\textgreek{’}}$à ce qui est tenu.

\subsubsection{Actions concrètes}
\begin{itemize}
\item Feedback positif structuré : partages où chacun repère ce qui fonctionne chez l$\text{\textgreek{’}}$autre.
\item Critiques constructives : formuler pour élever, pas pour soulager.
\item Suivi personnalisé : échanges réguliers (progrès, obstacles, objectifs).
\item Plans de développement coconstruits : objectifs mesurables + étapes + soutien.
\item Transparence : expliquer les décisions, partager les erreurs, assumer les arbitrages.
\item Engagement visible : être présent (cohérence et constance = crédibilité).
\end{itemize}
Dans une approche mindful engagement, la confiance n$\text{\textgreek{’}}$est pas seulement un résultat relationnel.
C$\text{\textgreek{’}}$est aussi un levier d$\text{\textgreek{’}}$apprentissage du leadership dont un plan de
développement applicable s$\text{\textgreek{’}}$appuie sur :

\begin{itemize}
\item Approche (besoins + état d$\text{\textgreek{’}}$esprit + identification des objectifs
d$\text{\textgreek{’}}$apprentissage)
\item Action (défis adaptés : hors zone de confort, mais pas destructeurs)
\item Réflexion (debrief, contrefactuels, leçons transférables) 
\end{itemize}
Sur le plan du développement, la confiance se renforce particulièrement lorsqu$\text{\textgreek{’}}$elle autorise des
défis réels. Des supports de leadership development rappellent que le développement repose sur un équilibre entre
challenge et support : trop peu de défi stagne, trop peu de soutien brise. Pour que l$\text{\textgreek{’}}$expérience
devienne un actif, encore faut-il la traiter activement. Ashford et DeRue montrent que la valeur développementale
d$\text{\textgreek{’}}$une expérience dépend moins de l$\text{\textgreek{’}}$événement en soi que de la manière dont il
est approché, vécu et ensuite réfléchi : l$\text{\textgreek{’}}$apprentissage exige une posture attentive et volontaire
\label{ref:RND3PBoy8YU0X}(Ashford \& DeRue, 2012). Dans la même logique, des pratiques de debrief après action doivent
être instaurées pour capter les leçons de l$\text{\textgreek{’}}$expérience par des after-event reviews. Pratiquer
\ des scénarios what-if, pour extraire des enseignements transférables. Ainsi encadrée, la confiance permet de vivre
des défis réels. Et les défis, bien digérés, renforcent la confiance.

\section{La performance comme conséquence mesurée d$\text{\textgreek{’}}$un système sain.}
Il faut mettre la performance en évidence, car cette mesure peut être comparée à soi (progression) ou à
l$\text{\textgreek{’}}$autre (classement). Cependant, la performance reste stérile si l$\text{\textgreek{’}}$être
n$\text{\textgreek{’}}$est pas libéré. Sans confiance, sa prise d$\text{\textgreek{’}}$initiative se contracte,
l$\text{\textgreek{’}}$erreur devient une menace, et l$\text{\textgreek{’}}$apprentissage se réduit.
L$\text{\textgreek{’}}$individu protège alors son ego, il évite l$\text{\textgreek{’}}$erreur et limite sa croissance.
Par cette logique, l$\text{\textgreek{’}}$expérience ne forme pas automatiquement. Elle doit donc être traitée avec
intention, expérimentation et réflexion \label{ref:RND52IQNNKBCF}(Ashford \& DeRue, 2012). Dès lors, la performance
durable devient moins une affaire de pression que de répétition, de sens et de régulation.

Ainsi, le maintien d$\text{\textgreek{’}}$un alignement éthique construit une performance plus robuste, car il refuse
les raccourcis toxiques. Les cadres déontologiques ne s$\text{\textgreek{’}}$excluent pas, ils se complètent pour
décider. Ils rappellent trois grandes visions structurantes (vertu, devoir, conséquences). Ils soulignent que les
dilemmes ne se résolvent pas seulement par l$\text{\textgreek{’}}$efficacité instrumentale, mais par la cohérence du
jugement. 

\subsubsection{Principes de leadership associés}
\begin{itemize}
\item Performance = apprentissage + répétition + sens (pas seulement pression).
\item Mesurer sans réduire (l$\text{\textgreek{’}}$indicateur sert le progrès, pas l$\text{\textgreek{’}}$identité).
\item Alignement éthique (la performance n$\text{\textgreek{’}}$a de valeur que si elle respecte les principes).
\end{itemize}
Dans ce contexte, un ethical leadership refuse l$\text{\textgreek{’}}$instrumentalisation, la tricherie ou
l$\text{\textgreek{’}}$humiliation, précisément parce que ces pratiques détruisent le respect, puis la confiance, et
finissent par dégrader la performance. La performance doit être rendue visible, sans pour autant réduire
l$\text{\textgreek{’}}$individu à un indicateur. Des objectifs explicites (de type SMART), des révisions et des
évaluations régulières, un journal de progression, et des rituels de reconnaissance matérialisent le chemin plutôt que
le seul résultat.

\subsubsection{Actions concrètes}
\begin{itemize}
\item Objectifs SMART : spécifiques, mesurables, atteignables, pertinents, temporels.
\item Révisions régulières : ajuster sans renier l$\text{\textgreek{’}}$ambition.
\item Évaluations périodiques : comparer aux objectifs, puis décider des correctifs.
\item Analyses comparatives : se comparer à soi-même (trajectoire) pour renforcer l$\text{\textgreek{’}}$agentivité.
\item Reconnaissance et célébration : ritualiser les progrès, pas seulement les podiums.
\item Journal / tableau de bord des réalisations : matérialiser la progression. 
\end{itemize}
Les études montrent que les comportements de leadership expliquent une part majeure de la variance des critères
d$\text{\textgreek{’}}$efficacité, et constituent un levier concret d$\text{\textgreek{’}}$amélioration. Dans leur
synthèse, la contribution relative des comportements apparaît particulièrement saillante : par exemple, pour
l$\text{\textgreek{’}}$efficacité du leader et la performance de groupe, la part de variance expliquée attribuée aux
comportements est rapportée comme très élevée \label{ref:RND3WG3brbyy3}(Derue et al., 2011). Plus le leadership devient
une pratique visible (cadre, feedback, soutien, recadrage, apprentissage post-action), plus la performance devient
probante.

\section{Adapter sans se trahir}
Un des pièges classiques du leadership consiste à se définir par un style unique. Or, l$\text{\textgreek{’}}$efficacité
se joue souvent dans la capacité à adapter les comportements à la situation. Les recherches montrent que les traits
influencent la probabilité de certains comportements, mais que ce sont surtout les comportements (orientés tâche,
relation, changement) qui expliquent une part majeure de l$\text{\textgreek{’}}$efficacité et
qu$\text{\textgreek{’}}$un modèle intégrant traits et comportements est particulièrement pertinent. Cette idée
s$\text{\textgreek{’}}$accorde avec les perspectives contemporaines qui mettent l$\text{\textgreek{’}}$accent sur la
contingence et sur la diversité des mécanismes de leadership selon le contexte \label{ref:RNDKRAcioXr55}(Dinh et al.,
2014). Elle s$\text{\textgreek{’}}$accorde aussi avec une lecture trait-based plus nuancée : les traits ne doivent pas
être considérés isolément, mais comme des constellations d$\text{\textgreek{’}}$attributs dont les effets peuvent être
distaux ou proximaux selon les processus en jeu \label{ref:RNDITsmHNuXCg}(Zaccaro, 2007).

Dans cette perspective, un point particulièrement utile concerne l$\text{\textgreek{’}}$émergence. Certaines personnes
plus réservées peuvent néanmoins émerger lorsqu$\text{\textgreek{’}}$elles adoptent ponctuellement les comportements
requis par la situation. Des résultats sur l$\text{\textgreek{’}}$extraversion montrent que le fait
$\text{\textgreek{«}}$d$\text{\textgreek{’}}$agir extraverti$\text{\textgreek{»}}$ peut constituer une cause proximale
de leadership émergent, y compris chez des individus moins extravertis au niveau dispositionnel
\label{ref:RNDMqlxd5aVFJ}(Spark \& O$\text{\textgreek{’}}$Connor, 2021). Cela n$\text{\textgreek{’}}$implique pas de
jouer un rôle en permanence, mais de mobiliser, quand nécessaire, un répertoire comportemental ciblé au service du
collectif, puis de revenir à un mode plus naturel pour préserver l$\text{\textgreek{’}}$équilibre.

\section{Conclusion}
Créer les conditions dans lesquelles les autres peuvent donner le meilleur d$\text{\textgreek{’}}$eux-mêmes sans se
perdre repose sur une architecture : le respect construit le cadre, la confiance ouvre
l$\text{\textgreek{’}}$engagement, et la performance devient une conséquence plutôt qu$\text{\textgreek{’}}$une
obsession. Le leadership se stabilise lorsque l$\text{\textgreek{’}}$éthique protège le cadre, lorsque
l$\text{\textgreek{’}}$expérience est transformée en apprentissage par la réflexion, et lorsque les comportements
observables et donc entraînables deviennent l$\text{\textgreek{’}}$interface entre intention et résultats.

Le leadership aligné peut se résumer ainsi :

\begin{itemize}
\item Éthique en amont (respect de règles tenables).
\end{itemize}
\begin{itemize}
\item Développement en dynamique (la confiance transforme l$\text{\textgreek{’}}$expérience en apprentissage).
\item Performance en aval d$\text{\textgreek{’}}$une progression durable (issue d$\text{\textgreek{’}}$un leadership
stable).
\end{itemize}
\subsection{}
\clearpage\section{Bibliographie}
Ashford, S. J., \& DeRue, D. S. (2012). Developing as a leader. Organizational Dynamics, 41(2), 146‑154.
https://doi.org/10.1016/j.orgdyn.2012.01.008

Day, D., Bastardoz, N., Bisbey, T., Reyes, D., \& Salas, E. (2021). Unlocking Human Potential through Leadership
Training \& Development Initiatives. Behavioral Science \& Policy, 7(1), 41‑54.
https://doi.org/10.1177/237946152100700105

Derue, D. S., Nahrgang, J. D., Wellman, N., \& Humphrey, S. E. (2011). Trait and Behavioral Theories of
\ Leadership\leavevmode\,: An integration and Meta‐Analytic test of their relative validity. Personnel Psychology,
64(1), 7‑52. https://doi.org/10.1111/j.1744-6570.2010.01201.x

Dinh, J. E., Lord, R. G., Gardner, W. L., Meuser, J. D., Liden, R. C., \& Hu, J. (2014). Leadership theory and research
in the new millennium\leavevmode\,: Current theoretical trends and changing perspectives. The Leadership Quarterly,
25(1), 36‑62. https://doi.org/10.1016/j.leaqua.2013.11.005

Dirks, K. T., \& Ferrin, D. L. (2002). Trust in leadership\leavevmode\,: Meta-analytic findings and implications for
research and practice. Journal of Applied Psychology, 87(4), 611‑628. https://doi.org/10.1037/0021-9010.87.4.611

Frazier, M. L., Fainshmidt, S., Klinger, R. L., Pezeshkan, A., \& Vracheva, V. (2017). Psychological
Safety\leavevmode\,: A Meta‐Analytic Review and Extension. Personnel Psychology, 70(1), 113‑165.
https://doi.org/10.1111/peps.12183

Judge, T. A., Piccolo, R. F., \& Kosalka, T. (2009). The bright and dark sides of leader traits\leavevmode\,: A review
and theoretical extension of the leader trait paradigm. The Leadership Quarterly, 20(6), 855‑875.
https://doi.org/10.1016/j.leaqua.2009.09.004

Kant, I., Gregor, M., \& Korsgaard, C. M. (2008). Groundwork for the metaphysics of morals. Cambridge University Press.

Koeslag-Kreunen, M., Van Den Bossche, P., Hoven, M., Van Der Klink, M., \& Gijselaers, W. (2018). When Leadership Powers
Team Learning\leavevmode\,: A Meta-Analysis. Small Group Research, 49(4), 475‑513.
https://doi.org/10.1177/1046496418764824

Legood, A., Van Der Werff, L., Lee, A., \& Den Hartog, D. (2021). A meta-analysis of the role of trust in the
leadership- performance relationship. European Journal of Work and Organizational Psychology, 30(1), 1‑22.
https://doi.org/10.1080/1359432X.2020.1819241

Spark, A., \& O$\text{\textgreek{’}}$Connor, P. J. (2021). State extraversion and emergent leadership\leavevmode\,: Do
introverts emerge as leaders when they act like extraverts? The Leadership Quarterly, 32(3), 101474.
https://doi.org/10.1016/j.leaqua.2020.101474

Wallace, D. M., Torres, E. M., \& Zaccaro, S. J. (2021). Just what do we think we are doing? Learning outcomes of leader
and leadership development. The Leadership Quarterly, 32(5), 101494. https://doi.org/10.1016/j.leaqua.2020.101494

Zaccaro, S. J. (2007). Trait-based perspectives of leadership. American Psychologist, 62(1), 6‑16.
https://doi.org/10.1037/0003-066X.62.1.6
\end{document}
